%
% ACL anthology style paper
% Trisattā: An Advaita Vedānta-Grounded Ontological Taxonomy for
% Hallucination Detection in Large Language Models
%
% Compile with:
%   pdflatex trisatta_paper.tex
%   bibtex trisatta_paper
%   pdflatex trisatta_paper.tex
%   pdflatex trisatta_paper.tex
%

\documentclass[11pt]{article}

% ── ACL / EMNLP style ─────────────────────────────────────────────────────────
\usepackage[hyperref]{acl2023}
\usepackage{times}
\usepackage{latexsym}
\usepackage[T1]{fontenc}
\usepackage[utf8]{inputenc}
\usepackage{microtype}
\usepackage{inconsolata}

% ── Additional packages ───────────────────────────────────────────────────────
\usepackage{amsmath, amssymb}
\usepackage{booktabs}           % professional tables
\usepackage{graphicx}
\usepackage{multirow}
\usepackage{xcolor}
\usepackage{url}
\usepackage{enumitem}
\usepackage{fontspec}           % for Devanagari/Unicode Sanskrit terms

% ── Colours ───────────────────────────────────────────────────────────────────
\definecolor{paramcol}{RGB}{34, 139, 34}   % forest green — pāramārthika
\definecolor{vyavcol}{RGB}{30, 100, 200}   % royal blue   — vyāvahārika
\definecolor{praticol}{RGB}{200, 50, 30}   % vermilion    — prātibhāsika
\definecolor{adhyasacol}{RGB}{150, 0, 150} % purple       — adhyāsa

% ── Macros ────────────────────────────────────────────────────────────────────
\newcommand{\param}{\textcolor{paramcol}{\textit{pāramārthika}}}
\newcommand{\vyav}{\textcolor{vyavcol}{\textit{vyāvahārika}}}
\newcommand{\prati}{\textcolor{praticol}{\textit{prātibhāsika}}}
\newcommand{\adhyasa}{\textcolor{adhyasacol}{\textit{adhyāsa}}}
\newcommand{\trisatta}{\textit{trisattā}}
\newcommand{\satta}{\textit{sattā}}
\newcommand{\vivarta}{\textit{vivartavāda}}
\newcommand{\system}{\textsc{Trisatt\={a}}}

% ── Metadata ──────────────────────────────────────────────────────────────────
\title{\system: An Advaita Ved\={a}nta-Grounded Ontological Taxonomy\\
       for Hallucination Detection in Large Language Models}

\author{
  Shrikantha Devaru \\
  \texttt{shrikanthdevaru@gmail.com}
}

\begin{document}
\maketitle

% ═══════════════════════════════════════════════════════════════════════════════
\begin{abstract}
We present \system{}, a three-tier hallucination taxonomy for Large Language
Models (LLMs) grounded in the \trisatta{} (three levels of being) doctrine of
Advaita Ved\={a}nta. Standard binary detection collapses epistemologically
distinct output types into a single undifferentiated class. Our framework maps
LLM outputs to three ontological levels: \param{} (factually verifiable),
\vyav{} (pragmatically plausible but unverified), and \prati{} (apparent
coherence that dissolves on inspection). We operationalise this taxonomy via
three complementary scores — embedding cosine similarity to a verified fact
store, NLI-based contextual entailment, and source-traceability — and introduce
an \adhyasa{} (superimposition) detector that isolates hallucinations disguised
as \vyav{} content. Evaluated on HaluEval ($n=100$) and TruthfulQA ($n=100$),
\system{} achieves 87\% and 75\% binary-equivalent accuracy respectively, while
characterising 11\%--26\% of items as \adhyasa{} cases invisible to binary
detection. We argue that \trisatta{} is \emph{strictly more expressive} than
binary hallucination detection and provides actionable, graded trust signals for
downstream systems. Code is available at: \url{https://github.com/shrikanthdevaru-lang/trisatta-hallucination}
\end{abstract}

% ═══════════════════════════════════════════════════════════════════════════════
\section{Introduction}

The problem of hallucination in Large Language Models --- the generation of
fluent, confident, but factually incorrect content --- has emerged as a critical
barrier to deployment in knowledge-sensitive domains \cite{ji2023survey}. Existing
detection frameworks overwhelmingly adopt a \textbf{binary paradigm}: an output
sentence is either hallucinated or not \cite{maynez2020faithfulness,lin2022truthfulqa}.

This binary framing, while computationally convenient, is epistemologically
impoverished. It conflates three fundamentally distinct output types:
\begin{enumerate}[leftmargin=*,noitemsep]
  \item \textbf{Verifiable outputs} --- claims for which truth-value is
    determinate and checkable against an external knowledge base;
  \item \textbf{Pragmatically coherent outputs} --- grammatically well-formed,
    semantically plausible, context-consistent, but externally unverified;
  \item \textbf{Apparently coherent outputs} --- claims whose surface coherence
    dissolves on factual inspection.
\end{enumerate}

We argue that Advaita Ved\={a}nta's \trisatta{} doctrine provides exactly the
right conceptual vocabulary for this three-way distinction, and that
operationalising it yields a classifier strictly more expressive than any binary
detector.

% ═══════════════════════════════════════════════════════════════════════════════
\section{Theoretical Framework}
\label{sec:theory}

\subsection{Advaita Vedānta Trisattā Doctrine}

Advaita Ved\={a}nta posits three levels of being (\satta{}):

\paragraph{\param{} \satta{} (absolute being).}
The level of ultimate, unchanging reality. We adapt this to designate claims
whose truth is \emph{absolutely grounded} in a verified, external knowledge base.

\paragraph{\vyav{} \satta{} (transactional / conventional being).}
The level of everyday empirical reality --- reliable within the phenomenal world
but not ultimate. We map this to claims that are \emph{pragmatically coherent}:
internally consistent, linguistically well-formed, contextually plausible, but
lacking direct factual verification.

\paragraph{\prati{} \satta{} (apparent / illusory being).}
The level of perceptual error --- the rope mistaken for a snake in dim light.
We map this to LLM outputs whose apparent coherence dissolves under verification.

\subsection{Adhyāsa and Vivartavāda}

\paragraph{\adhyasa{} (superimposition)} is the cognitive mechanism by which an
illusory object is mistaken for a real one. In our framework, \adhyasa{} fires
when \prati{} content passes surface-level coherence checks and is thereby
superimposed onto the \vyav{} level --- precisely the failure mode of binary
classifiers.

\paragraph{\vivarta{}} holds that apparent transformation occurs without
ontological change in the underlying substrate. An LLM hallucination is a
\textit{vivarta}: a statistically fluent output from a parametric distribution,
with no genuine epistemic grounding, yet appearing indistinguishable from
grounded output on the surface.

% ═══════════════════════════════════════════════════════════════════════════════
\section{Formal Classification Criteria}
\label{sec:formal}

Let $c$ be a claim (LLM output sentence) and $x$ be the context (prompt or
reference passage). We define three scores.

\paragraph{Score 1 --- Embedding Similarity ($s_1 \in [0,1]$).}
Cosine similarity between the embedding of $c$ and the nearest verified fact
in a FAISS-indexed knowledge base $\mathcal{K}$:
\begin{equation}
  s_1(c) = \max_{f \in \mathcal{K}}\, \cos\!\left(\mathbf{enc}(c),\, \mathbf{enc}(f)\right)
\end{equation}
This is the \param{} signal --- proximity to absolute factual ground truth.

\paragraph{Score 2 --- Internal Consistency ($s_2 \in [0,1]$).}
NLI entailment probability of $c$ given $x$:
\begin{equation}
  s_2(c, x) = P_\mathrm{NLI}(\mathrm{entailment} \mid \mathrm{premise}{=}x,\, \mathrm{hyp}{=}c)
\end{equation}
This is the \vyav{} signal --- pragmatic coherence within context.

\paragraph{Score 3 --- Source Traceability ($s_3 \in [0,1]$).}
Directed n-gram overlap between $c$ and $x$:
\begin{equation}
  s_3(c,x) = 0.6\cdot\frac{|\mathrm{ngrams}(c)\cap\mathrm{ngrams}(x)|}{|\mathrm{ngrams}(c)|}
             +0.25\cdot J(c,x) + 0.15\cdot\mathrm{PNO}(c,x)
\end{equation}
where $J$ is Jaccard overlap and $\mathrm{PNO}$ is proper-noun overlap.
Low traceability signals \prati{} dissolution.

\paragraph{Composite Score.}
\begin{equation}
  S(c,x) = 0.45\,s_1 + 0.30\,s_2 + 0.25\,s_3
\end{equation}

\paragraph{Decision Rule (priority order).}
\begin{align}
  s_1 \geq 0.82 \;\wedge\; s_2 \geq 0.80 &\;\Rightarrow\; \text{\param{}} \quad (L3) \\
  S \geq 0.55                              &\;\Rightarrow\; \text{\vyav{}}  \quad (L2) \\
  \text{otherwise}                         &\;\Rightarrow\; \text{\prati{}} \quad (L1)
\end{align}

\noindent\textbf{\adhyasa{} Signature:} fires when level $= L1$ and $0.40 \leq S < 0.65$.

% ═══════════════════════════════════════════════════════════════════════════════
\section{Implementation}
\label{sec:impl}

We implement \system{} as a modular Python package (\texttt{trisatta}) using:

\begin{itemize}[leftmargin=*,noitemsep]
  \item \texttt{all-MiniLM-L6-v2} \cite{reimers2019sentencebert} for $s_1$
  \item \texttt{cross-encoder/nli-MiniLM2-L6-H768} for $s_2$
  \item \texttt{faiss-cpu} \cite{johnson2019billion} for FAISS nearest-neighbour retrieval
\end{itemize}

The classifier returns \texttt{\{satta\_level, confidence, adhyasa\_signature, scores, explanation\}} for each (claim, context) pair. All thresholds are isolated in a single configuration module for reproducibility.

\paragraph{Key design note.}
For QA-style datasets where the ``context'' is a \emph{question} rather than a
declarative passage, NLI scores are systematically low (questions do not entail
their answers). We use the \texttt{knowledge} field in HaluEval (a declarative
supporting passage) as the NLI premise instead, which restores the full
discriminative power of Score 2.

% ═══════════════════════════════════════════════════════════════════════════════
\section{Experiments}
\label{sec:exp}

\subsection{Datasets}

\paragraph{HaluEval QA \cite{halueval2023}.}
100 items sampled from the QA hallucination evaluation benchmark. Each item
provides a knowledge passage, question, answer, and binary hallucination label.
We use the knowledge passage as the NLI context.

\paragraph{TruthfulQA \cite{lin2022truthfulqa}.}
100 items from the truthfulness benchmark. Correct answers form the
not-hallucinated class; incorrect answers form the hallucinated class, with
the question as context.

\subsection{Results}

\begin{table}[t]
\centering
\small
\begin{tabular}{lrr}
\toprule
\textbf{Trisattā Level} & \textbf{not-halluc.} & \textbf{halluc.} \\
\midrule
\param{} (L3) & \textbf{15} & 0 \\
\vyav{} (L2)  & 34 & 7 \\
\prati{} (L1) & 6  & \textbf{38} \\
\midrule
\textbf{Accuracy} & \multicolumn{2}{c}{\textbf{87\%}} \\
\textbf{Halluc. F1} & \multicolumn{2}{c}{\textbf{0.85}} \\
\textbf{\adhyasa{} rate} & \multicolumn{2}{c}{\textbf{11\%}} \\
\bottomrule
\end{tabular}
\caption{3$\times$2 confusion matrix on HaluEval ($n=100$). \param{} achieves zero false positives.}
\label{tab:halueval}
\end{table}

\begin{table}[t]
\centering
\small
\begin{tabular}{lrr}
\toprule
\textbf{Trisattā Level} & \textbf{not-halluc.} & \textbf{halluc.} \\
\midrule
\param{} (L3) & 2  & 0 \\
\vyav{} (L2)  & 22 & 1 \\
\prati{} (L1) & 24 & \textbf{51} \\
\midrule
\textbf{Accuracy} & \multicolumn{2}{c}{\textbf{75\%}} \\
\textbf{Halluc. F1} & \multicolumn{2}{c}{\textbf{0.80}} \\
\textbf{\adhyasa{} rate} & \multicolumn{2}{c}{\textbf{26\%}} \\
\bottomrule
\end{tabular}
\caption{3$\times$2 confusion matrix on TruthfulQA ($n=100$).}
\label{tab:tqa}
\end{table}

Tables~\ref{tab:halueval} and \ref{tab:tqa} show the full 3$\times$2 confusion
matrices. The \adhyasa{} rate is substantially higher on TruthfulQA (26\%)
than HaluEval (11\%), reflecting TruthfulQA's more epistemically ambiguous
questions where prātibhāsika content more readily appears vyāvahārika.

% ═══════════════════════════════════════════════════════════════════════════════
\section{Theoretical Contribution}
\label{sec:contribution}

Contemporary hallucination detection operates within a binary ontology: an
output sentence is classified as either \emph{hallucinated} or
\emph{not hallucinated} \cite{maynez2020faithfulness,ji2023survey}. While
computationally convenient, this dichotomy conflates epistemologically distinct
failure modes and obscures the generative mechanism that produces erroneous
outputs. We argue that Advaita Ved\={a}nta's \trisatta{} doctrine furnishes a
strictly more expressive and theoretically motivated taxonomy.

\paragraph{Expressiveness beyond binary detection.}
Binary classification partitions the output space into two equivalence classes.
The \trisatta{} taxonomy imposes a strict refinement: every binary class is
split into sub-classes with independent discriminating criteria.
\param{} content (L3) is verifiable against an external knowledge base --- its
truth-value admits empirical falsification. \vyav{} content (L2) is
pragmatically coherent but epistemically underdetermined. \prati{} content (L1)
exhibits apparent coherence that dissolves under verification. Binary classifiers
collapse L3 and L2 (both ``not hallucinated''), obscuring a critical epistemic
distinction: a \param{} claim is reliably trustworthy; a \vyav{} claim demands
independent verification.

\paragraph{\adhyasa{} as the mechanism of hallucination.}
The Vedāntic concept of \adhyasa{} identifies the precise cognitive error by
which an illusory object is mistaken for a real one --- the snake-rope
illusion. In LLM terms, \adhyasa{} occurs when \prati{} content passes
surface-level coherence filters and is thereby superimposed onto the \vyav{}
level. Our \adhyasa{} signature captures exactly this failure mode with a
measurable composite score, something no binary classifier can detect.

\paragraph{\vivarta{} and the apparent-real distinction.}
The doctrine of \vivarta{} holds that apparent transformation occurs without
ontological change. An LLM hallucination is a \textit{vivarta}: a statistically
fluent output from a parametric distribution that mimics but does not replicate
epistemic grounding. This motivates our three-score decomposition, each probing
a distinct ontological dimension that binary detectors treat as a single
undifferentiated feature.

\paragraph{Empirical consequences.}
The 3$\times$2 confusion matrix directly quantifies what binary systems cannot
express. On HaluEval, 7 of 45 hallucinated items land in the \vyav{} zone,
evading any binary detector. \param{} items achieve \textbf{zero false
positives} across both datasets --- a high-confidence trust tier binary systems
cannot represent. The \adhyasa{} rate provides a new diagnostic metric for
comparing LLM architectures beyond aggregate hallucination rates.

% ═══════════════════════════════════════════════════════════════════════════════
\section{Discussion}
\label{sec:disc}

\paragraph{Limitations.}
Threshold selection (\texttt{PARAM\_THRESH}, \texttt{VYAV\_THRESH}) requires
calibration per domain. The fact store must be pre-populated with verified
facts; its coverage directly limits \param{} recall. NLI scores are sensitive
to context type: questions produce systematically lower entailment scores than
declarative passages.

\paragraph{Future work.}
(1) Fine-grained \adhyasa{} sub-types (semantic vs.\ syntactic superimposition).
(2) Integration with RAG pipelines as a real-time grounding verifier.
(3) Extension to Sanskrit NLP corpora, where Vedāntic vocabulary has direct
ontological precision.
(4) Human evaluation of \trisatta{} labels for inter-annotator agreement.
(5) Fine-tuning a purpose-built NLI model on trisattā-annotated data.

% ═══════════════════════════════════════════════════════════════════════════════
\section{Related Work}
\label{sec:related}

Hallucination detection has been approached through faithfulness metrics
\cite{maynez2020faithfulness}, knowledge-grounded NLI \cite{honovich2022true},
and question-answering-based evaluation \cite{fabbri2021qafacteval}. More recent
work uses LLMs as self-evaluators \cite{manakul2023selfcheckgpt} or retrieval to
ground claims \cite{shi2023replug}. TruthfulQA \cite{lin2022truthfulqa} and
HaluEval \cite{halueval2023} provide standardised benchmarks. Our work differs
by providing a \emph{graded} rather than binary taxonomy, grounded in a formal
ontological framework rather than empirical heuristics.

The use of Indian philosophical frameworks in AI is nascent. \citet{ganeri2020}
has explored Ny\={a}ya epistemology for knowledge representation, and
\citet{bharati1996paninian} adapted P\={a}\d{n}inian grammar for NLP.
To our knowledge, this is the first work to operationalise \trisatta{} for
LLM evaluation.

% ═══════════════════════════════════════════════════════════════════════════════
\section{Conclusion}
\label{sec:conclusion}

We introduced \system{}, a novel three-tier hallucination taxonomy for LLMs
grounded in Advaita Ved\={a}nta's ontological framework. By operationalising
\param{}, \vyav{}, and \prati{} levels through embedding similarity, NLI
entailment, and source-traceability, we produce a classifier provably more
expressive than binary detection and reveal epistemologically critical
distinctions invisible to existing methods. Our \adhyasa{} detector isolates the
precise superimposition mechanism underlying LLM hallucination, providing a new
diagnostic metric for AI reliability research. We release the full
implementation and both evaluation datasets at
\url{https://github.com/shrikanthdevaru-lang/trisatta-hallucination}.

% ═══════════════════════════════════════════════════════════════════════════════
\section*{Acknowledgements}
[Acknowledgements to be added upon acceptance.]

% ═══════════════════════════════════════════════════════════════════════════════
\bibliography{trisatta_refs}
\bibliographystyle{acl_natbib}

\end{document}
